\documentclass[10pt,letterpaper]{article}
\usepackage[T1]{fontenc}
\usepackage{amssymb}
\usepackage{amsmath}
\usepackage[margin=0.5in]{geometry}
\usepackage{xcolor}
\title{Introduction to Quantum Electronic Devices and Applications Post Lecture 17 Exercise Problem 3}
\author{Jeremy Chen}
\date{30/03/2026}
\begin{document}
	\maketitle

\noindent We can estimate the critical temperature by solving for the temperature at which the lattice has enough energy available to break the coherent state. If $\Delta = 0.3$ meV, estimate the critical temperature $T_{C}$.

\color{violet}
With energy gap of $\frac{N}{2}\Delta$ and total lattice energy of $Nk_{B}T$, the critical temperature is
\begin{gather*}
\frac{N}{2}\Delta = Nk_{B}T \rightarrow T_{C} = \frac{0.15}{k_{B}}\ \text{meV} = 1.741 K
\end{gather*}
\color{black}

\end{document}